\section{Discussion}
\label{sec:discussion}

For the feature-based models, confidence that was close to calibrated on unseen artificial bearings did not carry over
to real damage: confidence thresholds, post-hoc calibration and conformal prediction all failed in the same direction,
towards overconfidence. Because every bearing belongs to exactly one split, the 5\,\% target is exceeded in all 105
rotations at the main condition, and even the lower ends of the bootstrap intervals over test bearings lie above it,
this is a property of the shift rather than of a particular choice of bearings. The 1D-CNN failed earlier, already on
unseen artificial bearings. The following subsections explain why confidence fails, what this means for calibration,
conformal prediction and practice, and where the study is limited.

\subsection{Why confidence fails on real damage}

The failures are not spread evenly over the real-damage bearings; they concentrate on four of them
(Section~\ref{sec:silent}). Their envelope spectra explain why. For the outer-race bearings KA15 and KA22, the peak at
the outer-race fault frequency is only 2.2 times the spectrum's noise floor (median over windows), and for the
inner-race bearing KI17 the peak at the inner-race frequency is also 2.2. Healthy bearings reach 1.5--2.6 and
1.6--2.1 on the same two features, whereas every artificially damaged bearing reaches at least 5.2 (outer race) or
3.8 (inner race). In the features the models see, these damaged bearings look healthy, and the models classify them
as healthy: 79--85\,\% of windows for KA15 and KA22, 58--70\,\% for KI17. KI16 shows a weak peak at the
\emph{outer}-race frequency (3.2) alongside a weak inner-race peak (2.8), and its windows are split between the two
fault classes.

This has two consequences. First, where the features separate the classes well, the dominant error is the most
dangerous one, a missed fault: at the main condition 62--68\,\% of all errors on real damage declare a damaged
bearing healthy, and 72--74\,\% of the errors made automatically at the 5\,\% target do; at low torque the share is
69--77\,\%. At lower radial force and lower speed, false alarms on healthy bearings are also common, and missed faults
account for only 20--27\,\% and 45--50\,\% of the errors (Section~\ref{sec:conditions}). Second, the models' confidence is not unreasonable given their inputs---a
spectrum without a fault peak \emph{is} typical of a healthy bearing in the training data. No post-hoc calibration
fitted on artificial damage can detect this, because artificial damage always produces a clear peak, and the
calibration set therefore never contains a damaged bearing that looks healthy. The shift is not a gradual drift of
the feature distribution that calibration could absorb; it introduces a type of case the source domain lacks.

The silence is not an artefact of one operating condition: KA15 and KA22 are misclassified at all four conditions,
KI17 at three and KI16 at two. Why the real damage is silent is a separate question that our data cannot settle. The documented
damage does not explain it: the silent bearings include fatigue pitting and an indentation (KA15), and KI16 has the
largest documented extent of all real bearings (level 3), whereas KA04, with the same damage mode and extent as KA22,
shows a clear peak (16.8 times the noise floor). One
explanation is that the damage produces weaker or less regular impacts than a sharp artificial defect; another is
that its impacts excite resonances outside the fixed 2--10\,kHz band. Selecting the band with the fast kurtogram did not
support the second explanation (Section~\ref{sec:adaptive}): the band it selected was dominated by impulsive content
present in healthy bearings as well, and the silent bearings stayed indistinguishable from healthy ones.

Our accuracy on real damage agrees with the dataset authors' own test: trained on artificially damaged bearings, their
classifiers reached 62--75\,\% on real damage with vibration features \citep{lessmeier2016}. Trained and tested on
different real-damage bearings, with time-, frequency- and wavelet-packet features, they reached up to 98\,\%,
recognising KA15 and KA22. The damage is therefore present in the vibration signal; what fails is the transfer of
features and confidence learned from artificial damage.

The 1D-CNN fails in a different way (Section~\ref{sec:cnnresults}). With the raw signal as input and only 12 to 14
training bearings, it partly recognises bearings by characteristics unrelated to their damage, so its confidence is
already misleading on unseen artificial bearings, and on real damage it raises false alarms far more often than it
misses faults. Unlike the failure of the feature-based models, which appears only after the shift, this problem is
visible in a validation by bearing. A random split of windows, which places windows of the same bearings in both
training and test sets, hides it completely: there the CNN is as accurate and as well calibrated as the feature-based
models (Section~\ref{sec:cnnresults}), in line with what \citet{vieira2026} found across several datasets.
Unsupervised domain adaptation did not repair the network's confidence either. AdaBN raised its accuracy on real
damage, but by moving its decisions rather than by telling it when they are wrong: it traded errors on outer-race and
healthy bearings for confident errors on inner-race bearings. Like other unsupervised adaptation methods, it aligns
the overall statistics of the two domains without any information on which target predictions are correct.

\subsection{Calibration in-domain does not imply calibration after the shift}

On unseen artificial bearings, the feature-based models are close to calibrated (ECE 0.01--0.09) and automate 81--91\,\% of cases
at or below 5\,\% error. On real damage, the same models are overconfident (ECE 0.17--0.20 after temperature scaling) and
can automate only 25--30\,\% of cases at that error even with a threshold chosen on the test data. A validation that
holds out artificial bearings---even one done carefully by bearing---would thus give a misleading picture of how the
system behaves in service.

The other operating conditions add a less intuitive observation: the better the features separate the classes on
artificial damage, the more misleading the confidence on real damage. At the main and low-torque conditions, where
in-domain accuracy is highest (0.81--0.94), the 5\,\% target is exceeded in every rotation. At lower radial force and
lower speed, in-domain accuracy is only 0.57--0.76, the models are less confident, and the error at the 5\,\% target
is lower (0.10--0.21), although still above the target in 68--99\,\% of the rotations that automate any case. A model that relies on a clear
fault signature confidently predicts ``healthy'' when that signature is missing, and a missing signature is exactly what
the silent bearings present. This mirrors findings on dataset shift in general machine learning \citep{ovadia2019},
here for a physically meaningful shift in a mechanical system.

\subsection{Conformal prediction under the shift}

The coverage guarantee of split conformal prediction assumes that calibration and test data are exchangeable, and the
artificial-to-real shift breaks this assumption. Coverage fell to 0.65--0.68 against a nominal 0.90, and the
prediction sets offered little warning: most contained a single class. Class-conditional conformal prediction did not
help (0.63--0.64). This differs from \citet{gonzalezgarcia2026}, where class-conditional thresholds partly repaired a
coverage failure caused by different class proportions between calibration and test recordings. Here the problem is
not the proportions but the content of a class: the real-damage bearings of a class do not resemble its artificial
calibration bearings. Weighted conformal prediction \citep{tibshirani2019} could in principle correct for covariate
shift, but it requires estimating how the test inputs differ from the calibration inputs, and the silent bearings
differ precisely by looking like another class.

\subsection{Implications for practice}

Our results suggest four recommendations for deploying automated bearing diagnosis.
\begin{enumerate}[nosep]
  \item \textbf{Do not set automation thresholds on artificial damage alone.} A threshold that gave 5\,\% error on
        held-out artificial bearings gave 24--27\,\% error on real damage.
  \item \textbf{Calibrate on labelled real bearings whenever possible, and expect to automate less.} Two real bearings
        per class reduced the error at the 5\,\% target from about 0.27 to 0.10--0.11, while the share of automated
        cases fell from 0.81--0.88 to 0.35--0.38. This is the honest trade-off: a system that knows less should act
        less.
  \item \textbf{Report results per bearing.} Averages over windows hid that the failures came from four bearings, and
        the wide bootstrap intervals show that a different set of test bearings could change the averages
        considerably.
  \item \textbf{Check features and models on unseen bearings.} Amplitude-based features reflected individual
        bearings, including healthy ones, and lowered accuracy on unseen artificial bearings from 0.88--0.91 to
        0.49--0.77; the 1D-CNN reached only 0.46. On real damage, however, no feature set or model we tested gave
        trustworthy confidence, including the network adapted to unlabelled real-damage data.
\end{enumerate}

\subsection{Limitations}
\label{sec:limitations}

The most important limitation is generality. All bearings come from one test rig and are of one type, the real
damage was produced in accelerated lifetime tests and may differ from damage that develops in the field, and bearings
with combined inner- and outer-race damage were excluded; the size of the effects may therefore differ on other
machines. Repeating the evaluation at all four operating conditions of the dataset shows that the failure is not tied
to one condition, but the experiment with labelled real bearings for calibration, the feature ablation, the adaptive
band and the 1D-CNN were run at the main condition only.

The number of bearings is small---11 with real damage and 6 healthy---so the bootstrap intervals over bearings are
wide. The main conclusions therefore rest on effects that hold in every rotation at the main condition, such as the
5\,\% target being exceeded in all of them, rather than on precise values. Each model was trained with a single random
seed; the rotations and the bootstrap capture variation over bearings, not over training randomness, which matters
most for the 1D-CNN. Conformal calibration is performed on windows, and windows from the same bearing are correlated,
so even the in-domain exchangeability assumption holds only approximately.

Several alternatives remain untested. The features use a fixed demodulation band and window length; the kurtogram
band we tested favoured impulsive content of the test rig, and other criteria for choosing the band were not
explored. The deep-learning baseline is a single architecture trained from scratch with fixed settings, and AdaBN was
the only domain adaptation tested. Pre-training on other data, stronger data augmentation, or adversarial and
discrepancy-based adaptation such as DANN \citep{ganin2016} or Deep CORAL \citep{sun2016} might generalise better
across bearings, although, like AdaBN, they use no labels from the target domain. Whether such methods make the
confidence trustworthy, and not only the accuracy, remains to be tested.
